\section{Related Work}
\label{sec:related_work}

Our research intersects four key literature domains: long-horizon error compounding, process verification in language agents, execution recovery and rollback, and verification scheduling.

\subsection{Error Compounding in Long-Horizon Autonomous Workflows}
Extensive empirical analyses demonstrate that autonomous agent performance degrades rapidly as workflow horizon expands~\citep{sinha2025illusion,wang2026reasoning}. In multi-step execution graphs, unverified errors compound: an invalid schema mutation or corrupted entity key in an early step causes subsequent transformations to operate on poisoned data, creating cascading failure cascades~\citep{chen2026horizon}. In sequential architectures, task success probability decays exponentially with horizon length $H$, formalized as $S(H) \approx \prod_{t=1}^H p_t$, where non-stationary environmental shifts further degrade per-step success $p_t$~\citep{wang2026reasoning}. While search algorithms like Tree of Thoughts (ToT)~\citep{yao2023tot} and Language Agent Tree Search (LATS)~\citep{zhou2024lats} explore alternative reasoning paths in closed puzzle environments, they do not address long-horizon execution workflows with concrete dataflow artifacts and persistent state dependencies.

\subsection{Process Verification in Agentic Systems}
To detect errors before trajectory completion, recent architectures incorporate process-level verification~\citep{venkataramani2026masprove,afoud2026vreact}. MAS-ProVe~\citep{venkataramani2026masprove} examined process-reward models and language-model-as-judge paradigms across multi-agent systems, finding that uncalibrated model-based verification suffers from high variance and inconsistent performance improvements. 

Addressing this variance, VReAct~\citep{afoud2026vreact} proposed a formal V-model architectural framework separating requirement decomposition from multi-tier verification gates. Crucially, VReAct introduced the concept of an \emph{integration-level verifier} designed to validate whether outputs from parallel dependency branches structurally compose before downstream execution. However, the pilot study in~\citep{afoud2026vreact} evaluated only unit-level and acceptance-level verifiers on a small 47-run question-answering dataset, explicitly omitting an empirical evaluation of the integration verifier and leaving verification scheduling unexamined. JevAct~\citep{shoham2026jevact} proposed action-level decision gates, but evaluated local tool call filtering rather than workflow-wide scheduling under global call constraints.

\subsection{Checkpoints, Rollback, and State Recovery}
Once an error is detected, an agent must recover without repeating the entire trajectory from scratch. Early reflection frameworks such as Reflexion~\citep{shinn2023reflexion} relied on verbal memory refinement, which struggles in multi-branch workflows with persistent state side-effects. 

Recent state-preservation frameworks introduce concrete environment checkpointing. ScienceFlow~\citep{zhao2026scienceflow} introduced the ESTRA execution controller, which checkpoints workspace files, memory, and validation logs, restoring historical states when validation evidence falls below progress thresholds. LEAD~\citep{pushkin2026lead} demonstrated that lookahead rollouts with checkpoint filtering break the ``no-recovery bottleneck'' in long-horizon reasoning. Similarly, AgentRewind~\citep{zhuang2026agentrewind} and RIR~\citep{yu2026rir} explored aligned environment-context rollback, investigating how much execution context must be restored to resume execution successfully. While these systems demonstrate the efficacy of rollback recovery, they focus on state reconstruction mechanics rather than the mathematical optimization of \emph{when} checkpoints should be tested.

\subsection{Verification Scheduling and Frequency}
Despite widespread adoption of intermediate checking, the temporal placement of verification checks remains poorly understood. Agentic Robot~\citep{yang2025agenticrobot} provided a rare empirical study on verification frequency, comparing fixed periodic intervals of 10, 20, and 50 execution steps in an embodied visual navigation task. Their results indicated that increasing frequency from 50 to 20 steps significantly improved completion, while increasing frequency from 20 to 10 steps yielded diminishing returns while doubling runtime latency. 

However, prior work has not compared fixed-periodic cadences against topological, dependency-aware milestone scheduling under an \textbf{identically matched verification budget}. Our work bridges this exact gap: holding the verifier, task graphs, fault injections, and available check budgets strictly constant, we directly test whether graph-derived dependency integration boundaries offer superior error containment compared to clock-like periodic intervals.
